\documentclass[a4paper,12pt]{article}

\usepackage[utf8]{inputenc}
\usepackage[T1]{fontenc}
\usepackage[brazil]{babel}
\usepackage{geometry}
\usepackage{longtable}
\usepackage{array}
\usepackage{enumitem}

\geometry{margin=2.5cm}

\title{Escopo – Subgrupo de Software}
\date{}

\begin{document}

\maketitle

\section{Frentes de Trabalho}

\subsection{Software Embarcado}

\begin{itemize}
    \item Leitura e interpretação dos sensores de distância (detectar paredes nas células adjacentes);
    \item Monitoramento da posição atual no labirinto (odometria e controle de orientação);
    \item Construção incremental do mapa do labirinto em tempo de execução;
    \item Tomada de decisão de rota com algoritmo de exploração (Flood Fill ou similar);
    \item Detecção da chegada na área objetivo (sala central do labirinto);
    \item Transmissão dos dados de telemetria para o sistema web via comunicação sem fio.
\end{itemize}

\subsection{Sistema Web de Telemetria}

\begin{itemize}
    \item Receber dados em tempo real via WebSocket ou protocolo MQTT;
    \item Exibir em tempo real os seguintes dados de telemetria:
    \begin{itemize}
        \item Tipo do labirinto (4×4, 8×8 ou 16×16);
        \item Trajeto percorrido, visualizado sobre o mapa do labirinto;
        \item Consumo de bateria;
        \item Velocidade média;
        \item Tempo de conclusão;
        \item Desafio cumprido (Sim / Não).
    \end{itemize}
    \item Persistir os dados finais em banco de dados ao término de cada corrida;
    \item Permitir consultas históricas por labirinto individual ou por todos os labirintos;
    \item Exibir os dados armazenados em tela de histórico com filtros.
\end{itemize}

\section{Entregáveis por Marco}

\begin{longtable}{|p{4cm}|p{7cm}|p{2cm}|}
\hline
\textbf{Entregável} & \textbf{Descrição} & \textbf{Marco} \\
\hline
Algoritmo de exploração (embarcado) & Código funcional de navegação com Flood Fill & PC2 \\
\hline
Módulo de sensoriamento & Leitura e interpretação dos sensores de distância & PC2 \\
\hline
Módulo de odometria & Monitoramento de posição e orientação do robô & PC2 \\
\hline
Transmissor de telemetria & Envio dos dados do embarcado para o servidor web & PC2 \\
\hline
Backend da aplicação web & API que recebe, processa e persiste os dados & PC2 \\
\hline
Frontend de telemetria & Interface web com dados ao vivo durante a corrida & IPP \\
\hline
Banco de dados e histórico & Persistência e consulta dos dados finais & IPP \\
\hline
Sistema integrado e testado & Todos os módulos integrados nos três labirintos & APT \\
\hline
Documentação técnica (UML) & Diagramas de casos de uso, classes, sequência e arquitetura & Por entrega \\
\hline
\end{longtable}

\section{Fora do Escopo}

Os itens abaixo não são responsabilidade do subgrupo de Software:

\begin{itemize}
    \item Projeto e montagem da estrutura física do micromouse (Subgrupo de Estrutura);
    \item Seleção, aquisição e conexão dos componentes eletrônicos (Subgrupo de Hardware);
    \item Cálculo e gerenciamento do consumo energético e da bateria (Subgrupo de Energia);
    \item Construção física do labirinto de testes 4×4 (responsabilidade coletiva do grupo).
\end{itemize}

\section{Premissas e Dependências}

\begin{itemize}
    \item O microcontrolador escolhido pelo Subgrupo de Hardware suporta comunicação sem fio e permite programação em Python;
    \item O Subgrupo de Hardware fornecerá interface de acesso aos sensores e encoders com documentação mínima;
    \item A apresentação final ocorrerá em ambiente com acesso a rede local (Wi-Fi);
    \item O grupo terá acesso a pelo menos um labirinto 4×4 para testes de integração antes do PC2.
\end{itemize}

\section{Restrições}

\begin{itemize}
    \item Nenhuma solução de software pronta pode substituir o desenvolvimento próprio;
    \item Ferramentas de inteligência artificial generativa não podem ser usadas nas atividades avaliativas;
    \item O código-fonte e a memória do labirinto não podem ser alterados durante a resolução de um trajeto;
    \item Cada issue no GitHub deve ter no máximo dois responsáveis.
\end{itemize}

\section{Critérios de Aceitação}

\begin{longtable}{|p{7cm}|p{6cm}|}
\hline
\textbf{Critério} & \textbf{Como será verificado} \\
\hline
Micromouse navega autonomamente no labirinto 4×4 & Demonstração na Inspeção Prévia (IPP) \\
\hline
Micromouse navega nos labirintos 8×8 e 16×16 & Demonstração na Apresentação Final (APT) \\
\hline
Dados de telemetria aparecem em tempo real no site & Observação direta durante a demonstração \\
\hline
Dados são armazenados no banco após cada corrida & Consulta na interface de histórico \\
\hline
Interface permite filtrar por labirinto ou por todos & Teste funcional na tela de histórico \\
\hline
Diagramas UML coerentes com o sistema implementado & Revisão no relatório técnico por Ponto de Controle \\
\hline
\end{longtable}

\section{Tecnologias Previstas}

\begin{table}[h!]
\centering
\begin{tabular}{|p{3.5cm}|p{4cm}|p{6.5cm}|}
\hline
\textbf{Camada} & \textbf{Tecnologia} & \textbf{Justificativa} \\
\hline
Embarcado & Python libESP\_IDF / C/C++ & Padrão para microcontroladores com suporte a bibliotecas de hardware \\
\hline
Comunicação & Wi-Fi + MQTT ou WebSocket & Protocolo leve para transmissão de dados em tempo real \\
\hline
Backend & Python (FastAPI) ou Node.js & Frameworks leves para APIs REST e WebSocket \\
\hline
Banco de dados & SQLite ou PostgreSQL & SQLite para testes locais; PostgreSQL para maior robustez \\
\hline
Frontend & HTML + JS ou React & Interface web acessível por navegador sem instalação \\
\hline
Versionamento & Git + GitHub & Exigido pelo Plano de Ensino com issues e sub-issues \\
\hline
Modelagem & Draw.io / Lucidchart & Criação dos diagramas UML exigidos nos Pontos de Controle \\
\hline
\end{tabular}
\end{table}

\end{document}